\section{Data and SDE-based Generation}
\label{sec:data_sde}
%ODE-only simulation vs NIH mismatch
%drift=OGTT, diffusion=데이터 기반 → 생성 데이터 품질 및 성능 향상
%단, 구조적 한계는 유지됨을 명확히

\subsection{SDE-based data generation for OGTT}
\label{subsec:sde}

For OGTT, we found a non-negligible gap between deterministic simulations of the mechanistic ODE model and patient trajectories.
We model this discrepancy by augmenting the ODE with a stochastic component and generate synthetic training data from a calibrated SDE.

\subsubsection{From deterministic mismatch to a bias-corrected SDE}
\label{subsubsec:sde-form}

Let $X(t)=(G(t),I(t),N_5(t),N_6(t))$ denote the OGTT state variables and $\theta=(S_I,\sigma)$ the parameters.
We take the OGTT ODE as the drift term and consider an SDE of the form
\begin{equation}
\label{eq:sde-form}
dX_t = \bigl[f(X_t,\theta) + \mu_{\mathrm{bias}}(t)\bigr]\,dt + \Sigma_{\mathrm{diff}}(t)\,dW_t,
\end{equation}
where $W_t$ is a Wiener process.
Here, $\mu_{\mathrm{bias}}(t)$ is a time-dependent drift correction accounting for systematic model bias,
and $\Sigma_{\mathrm{diff}}(t)$ is a time-dependent diffusion term capturing stochastic variability.
This formulation was motivated by the empirical observation that deterministic residuals are not mean-zero, suggesting a drift mismatch rather than purely random noise.

\begin{figure}[t]
    \centering
    % \includegraphics[width=0.95\linewidth]{figures/sde_effect_visualization.png}
    \fbox{\parbox{0.95\linewidth}{\vspace{2.6cm}\centering
    \textit{Placeholder: ODE vs SDE ensemble vs NIH trajectories (OGTT).}\vspace{2.6cm}}}
    \caption{OGTT trajectories: deterministic ODE versus calibrated SDE ensemble.
    The dashed curve shows the deterministic ODE simulation, and the black curve shows a representative NIH trajectory.
    While the mean of the SDE ensemble remains close to the ODE trajectory, the stochastic spread induced by the diffusion term produces a band
    that more realistically covers the variability observed in the NIH data.}
    \label{fig:sde-effect}
\end{figure}

\subsubsection{Calibration: diffusion from increment statistics (quadratic variation)}
\label{subsubsec:sde-diffusion}

A key point in calibrating \eqref{eq:sde-form} is that the diffusion term represents uncertainty in \emph{increments} rather than accumulated state residuals.
Using the variance of state residuals $\mathrm{Var}(Y_{\mathrm{real}}(t)-Y_{\mathrm{sim}}(t))$ leads to a dimensional inconsistency and can
overestimate diffusion, producing unrealistically large random-walk behavior.

Instead, we estimate diffusion from increment residuals over each observation interval.
Let $t_0<\cdots<t_{N-1}$ be the OGTT sampling times and define increments $\Delta Y_i := Y(t_{i+1})-Y(t_i)$.
Denoting the increment residual by $\Delta R_i := \Delta Y^{\mathrm{real}}_i - \Delta Y^{\mathrm{sim}}_i$,
the quadratic-variation relation suggests
\begin{equation}
\label{eq:diffusion-increment}
\mathrm{Var}(\Delta R_i)\approx \sigma_Y(t_i)^2\,\Delta t_i,
\qquad
\sigma_Y(t_i) \approx \sqrt{\frac{\mathrm{Var}(\Delta R_i)}{\Delta t_i}},
\qquad Y\in\{G,I\},
\end{equation}
with $\Delta t_i:=t_{i+1}-t_i$.
In practice, we compute interval-wise $\sigma_G(t_i)$ and $\sigma_I(t_i)$ after robust outlier removal.

\subsubsection{Calibration: drift bias from mean increment residuals}
\label{subsubsec:sde-bias}

To correct systematic mismatch, we estimate a drift-bias term from the mean increment residual:
\begin{equation}
\label{eq:drift-bias}
\mu_{\mathrm{bias},Y}(t_i) \approx \frac{\mathbb{E}[\Delta R_i]}{\Delta t_i},
\qquad Y\in\{G,I\}.
\end{equation}
Empirically, the mean increment residual is not negligible across time intervals, supporting the use of a drift correction in \eqref{eq:sde-form}.

\subsubsection{Estimated drift and diffusion parameters}
\label{subsubsec:sde-estimates}

Figure~\ref{fig:sde-dist-analysis} visualizes the interval-wise increment statistics used to calibrate $\mu_{\mathrm{bias}}(t)$ and $\sigma_Y(t)$
for glucose and insulin.
Table~\ref{tab:sde-calib} reports the corresponding numerical estimates (computed on the training subset).

\begin{figure}[t]
    \centering
    % \includegraphics[width=0.95\linewidth]{figures/sde_distribution_analysis.png}
    \fbox{\parbox{0.95\linewidth}{\vspace{2.6cm}\centering
    \textit{Placeholder: increment-residual distribution and state-dependence analysis (OGTT).}\vspace{2.6cm}}}
    \caption{Time-wise distributional calibration for OGTT SDE generation.
    For each observation interval, we analyze the distribution of increment residuals $\Delta R$,
    potential state dependence, and variability range, supporting the increment-based calibration in
    \eqref{eq:diffusion-increment}--\eqref{eq:drift-bias}.}
    \label{fig:sde-dist-analysis}
\end{figure}

\begin{table}[t]
\centering
\caption{Calibrated drift and diffusion statistics for OGTT SDE generation (train set only, $n=415$).
Values are estimated interval-wise from increment residuals; see \eqref{eq:diffusion-increment}--\eqref{eq:drift-bias}.}
\label{tab:sde-calib}
\begin{tabular}{rcccc}
\hline
$t$ (min) & $\mu_G(t)$ & $\sigma_G(t)$ & $\mu_I(t)$ & $\sigma_I(t)$ \\
\hline
0   & $-0.5638$ & $2.8235$ & $-0.1771$ & $4.5862$ \\
30  & $0.5693$  & $4.1121$ & $-0.1012$ & $6.7049$ \\
60  & $0.0438$  & $3.4084$ & $0.3671$  & $6.6500$ \\
90  & $0.0284$  & $3.2451$ & $0.3701$  & $5.6393$ \\
120 & $0.0284$  & $3.2451$ & $0.3701$  & $5.6393$ \\
\hline
\end{tabular}
\end{table}

For convenience, the entry at $t=120$ repeats the last interval estimate.
We additionally enforce physical clamping bounds based on NIH maxima:
$G_{\max}=506.4$ and $I_{\max}=788.64$.

\subsubsection{Simulation procedure}
\label{subsubsec:sde-sim}

We simulate \eqref{eq:sde-form} using an Euler--Maruyama scheme on a fine time grid, followed by subsampling to the OGTT observation times.
The calibrated time-dependent coefficients in Table~\ref{tab:sde-calib} are treated as piecewise-constant over each observation interval.
Implementation details are provided in the accompanying codebase.

Finally, we emphasize that the SDE-based data generation improves realism of training data by capturing variability around the deterministic drift.
However, it does not resolve structural limitations of glucose-only parameter recovery, which are analyzed in Section~\ref{sec:correctness}.